% RESEARCH_DRAFT_SHA256: 9d3019d5e1d66f382217f81aa59c13dc6bd94fc7f1bb4a2be29c5b3379a766f2
% Layout and source-reference wrappers expanded for the web renderer.
\documentclass[11pt]{article}
\usepackage[margin=25mm]{geometry}
\usepackage{amsthm}
\usepackage{fontspec,unicode-math}
\setmainfont{Latin Modern Roman}
\setmathfont{Latin Modern Math}
\usepackage{amsmath,amssymb,mathtools,xurl,hyperref}
\hypersetup{colorlinks=true,urlcolor=blue,linkcolor=blue}
\setcounter{secnumdepth}{0}
\setlength{\parindent}{0pt}
\setlength{\parskip}{5pt}
\setlength{\emergencystretch}{3em}
\newcommand{\sref}[2]{\hyperlink{source:#1:#2}{[S#2]}}
\newcommand{\cataloguescope}[1]{\paragraph{Source record.}#1}
\newcommand{\floor}[1]{\left\lfloor#1\right\rfloor}
\newtheorem{proposition}{Proposition}
\newtheorem{lemma}[proposition]{Lemma}
\begin{document}
\paragraph{Research manuscript.} OpenAI \texttt{gpt-6-astra}, reasoning \texttt{ultra}; 11 October 2026.
\paragraph{Preserved source body.} The complete catalogue mathematical text and sources follow with layout and source links expanded, including source qualifications and its historical preparation notes. New research begins at the section headed Research attempt.

\section{A typed plural-paper certificate for mathematical communication}\label{statement}

\subsection{Definitions and mathematical statement}

\paragraph{Shared notation.}Standard number systems are $\mathbb N,\mathbb Z,\mathbb Q,\mathbb R,\mathbb C$. Specialized notation and index ranges follow the cited source and the definitions below.


Mathematical communication presents definitions, results, arguments, or explanatory models to an audience. A communication medium is the format through which those are shared. Evaluation depends on the intended audience and objective, such as understanding an argument, finding a resource, or accessing a presentation.
\paragraph{Record scope and qualifications.}
The catalogue title is retained as metadata. Named certificates, contracts, bounds or frameworks in that title are not asserted as results of the original source extract.
The source's reference to alternatives in 2023 is historical. No current completeness claim, selected audience, or effectiveness metric is asserted.

Catalogue screening limitation: Workshop planning or software infrastructure agenda; the supplied passage does not identify a determinate mathematical target.

The source snapshot labels this record partially solved; this is catalogue metadata rather than independent verification of a mathematical result or the current completion of a project.
\paragraph{Statement recorded in the supplied source.}
The goals of this group are to reevaluate the aims of math communication, using specific examples, and introduce new concepts that they believe will impact the future of how mathematics is communicated across a wide audience.
In addition to thinking about how mathematics might be communicated in the future, this group also aims to outline existing alternatives that already exist in 2023. The concrete product of such a project would be a collection of resources which enable people to find alternative (and potentially) better ways of sharing their work. \href{http://aimpl.org/cyberinfrastructure/5/}{[S2]}
\paragraph{Source status.}
The supplied catalogue labels this source record partially solved; that metadata was checked against the supplied snapshot on 2026-10-11. The label is preserved for provenance and is not an independent verification of present open status or mathematical correctness.
\subsection{Short English statement}
Review mathematical communication practices and collect alternative ways to share mathematical work.
\subsection{Sources}
\begin{itemize}
\item[S1] Ulam AI and the UnsolvedMath Contributors, UnsolvedMath: A Curated Collection of Open Mathematics Problems, catalogue entry. Catalogue curation and metadata: CC BY 4.0. The source extract is retained; definitions, English summaries and qualifications are editorial additions.
\url{https://huggingface.co/datasets/ulamai/UnsolvedMath}.
\item[S2] AIM Workshop Problem Lists, source linked by the catalogue.
\url{http://aimpl.org/cyberinfrastructure/5/}.
\end{itemize}
\paragraph{Preparation and review.}
The original source and UnsolvedMath contributors retain attribution for the source extract. Codex (GPT-6; exact serving revision not exposed) assisted with definitions, English summaries, scope notes and formatting. No new human mathematical authorship or independently verified solution is asserted. Catalogue identity checked against the complete supplied registry and current repository. Distinct upstream records, including repeated questions and broad or damaged extracts, are retained for complete coverage of the source CS category. A repeated record is not a new independent conjecture. Complete mathematical-equivalence and current-literature review has not been performed. Preparation and final review follow contribution prompts C and D, with source gaps disclosed. Structural and rendering checks do not establish a new theorem or open status.
\label{end:statement}

\section{Research attempt}
\subsection{Literature, source context and scope}
The complete catalogue body above retains the original agenda, source status, and historical preparation. The AIMPL problem page could not be retrieved. The supplied complete source transcription was read, and the accessible primary AIM workshop announcement was checked for context. The announcement is not a substitute for verification of the inaccessible problem page. The catalogue's earlier proposed frameworks and research summaries are context, not established solutions credited to this attempt. The following is an explicitly limited investigation of one route; it neither substitutes a formalized toy model for the entire source agenda nor claims empirical or deployed-software results. Elementary auxiliary facts are not new original-target progress.

\subsection{Testing whether compact claims can discard their assumptions}
The communication agenda is broader than a mathematical theorem. We examine one proposed simplification: communicate an implication while suppressing its domain. Consider the same displayed sentence, ``$x^2=1$ implies $x=1$,'' attached either to the domain $D_+=\{x\in\mathbb R:x>0\}$ or to $D=\mathbb R$.
\begin{proposition}
The implication holds for every $x\in D_+$ but fails for $D$. Any decoder receiving only the displayed sentence cannot recover a correct yes/no universal-validity judgment for both original domain-bearing records.
\end{proposition}
\begin{proof}
Over the reals, $x^2-1=(x-1)(x+1)$. If this product is zero, one factor is zero, so $x=1$ or $x=-1$. Positivity excludes the latter and proves the implication over $D_+$. Over $D$, $x=-1$ is a counterexample. Erasing the domain maps the two records to the identical message. A decoder of that message has only one output, whereas correct validity judgments differ, so it cannot be correct on both.
\end{proof}
This blocks the attempted inference that identical rendered assertions have identical mathematical content independently of audience context. It suggests carrying explicit assumptions with a claim when context cannot supply them; it does not prescribe a universal communication medium or show that every omitted assumption causes confusion.
\subsection{Verification and remaining gaps}
The premise and conclusion were checked exactly at $x=-1,1$, producing the counterexample and the valid positive root. No resource collection, comparative communication study, or new community communication framework was produced. The elementary information-loss example addresses one small constraint on the source's much broader agenda.

\section{References and source attribution}
Primary source: \url{http://aimpl.org/cyberinfrastructure/5/}. Attempted primary AIMPL page retrieval; failed with timeout or gateway error. The supplied full source transcription was read; the inaccessible live page was not independently inspected.

\url{https://aimath.org/pastworkshops/cyberinfrastructure.html}. Read the primary December 2023 workshop announcement; it describes the open-source mathematical infrastructure agenda.

All retrievals were checked on 11 October 2026. Original contributor and catalogue attribution remains in the preserved Sources section. The elementary arguments below are self-contained; their exposition is not a claim of first discovery.


\paragraph{Integration review.} Codex (GPT-6; exact serving revision not exposed) reviewed the written arguments and expanded layout and source-reference wrappers. Repository metadata and the production web renderer were checked using Python 3.13 and JavaScript with MathJax 3.2.2. This model-assisted review is not independent human proof verification; no PDF compilation or formal proof checking is asserted.

\section{Completion Estimate}
\noindent\textbf{COMPLETION: 0\%}
No substantiated new progress toward the entire original target is claimed. The worked results and failed route are delimited above; a known result, illustrative model, or auxiliary obstruction is not counted as original-target completion.

\section{Final status and attribution}
\textbf{Attempt status: unresolved.} This describes this research attempt, not a new ruling on the source catalogue's status or on external literature claims.
The mathematics and this self-audit were generated by OpenAI \texttt{gpt-6-astra}, reasoning \texttt{ultra}, on 11 October 2026. Codex (GPT-6; exact serving revision not exposed) orchestrated the batch and contributed formatting/repository review. No human mathematical authorship, independent human proof checking, or formal verification is claimed. The historical preparation paragraphs in the reproduced catalogue body concern that body only.

Contribution prompt A was applied to the complete definition and supplied source record, followed by the adversarial checks of prompt B. Local reproduction files are \texttt{batch013/checks.py}, \texttt{computations.json}, and the per-ID audit JSON. They were not public links at preparation time. Numerical checks are finite and do not prove unrestricted results. No LaTeX engine was found locally; compilation was not performed by this research agent.

\end{document}
